\section{Watcom}

\begin{wrapfigure}[9]{r}{0.25\textwidth}
\centering
\includegraphics[width=.25\textwidth]{drawings/watcom.pdf}
\end{wrapfigure}


The DOS extender was magical but it was hard to setup as a standalone product. A bootstrap which would locate \cw{DOS/4GW}, the program and setup both required multiple steps and close to 100 lines of C code\footnote{Source: Watcom C/C++ Programmer's Guide, "7.1.1 The Stub Program"}. If there was no impact on the program itself, the build process was affected. The ramp up time was significant and raise the entry barrier. What was really needed was an integrated environment where the compiler and the linker would take care of bundling the extender and the application together into one executable.\\ 
\par

In 1993 the best compiler available on MS-DOS was irrevocably Watcom C 9.0\footnote{Editor's choice -- PC Magazine, April 1995} \footnote{PC Magazine ran a comparative study of a collection of C compilers.  They hailed WATCOM C/C++ v9.5 for producing the fastest code (by a *wide* margin).}. Started in 1986 when a team of undergraduate students at the University of Waterloo in Canada developed a FORTRAN compiler. The small team of three\footnote{The company success quickly grew the compiler team to 8 peoples.} made of Fred Crigger, Ian McPhee, and Jack Schueler improved the compiler to support multiple language and platforms. Their had chosen a lightning for their logo to advertise their focus on performances. On BSS you could ran into these kind of claims:\\
\par
\fq{WATCOM C/C++ will produce code which is at *least* twice as fast as your current 16 bit compiler, and more typically around five times as fast.}{rec.games.programmer}
\\
You could not open a computer magazine without finding a full page advertising Watcom's compiler. All ads ostensibly underlined the presence of a Dos Extender.
\par
\fullimage{watcom_ad.png}
\par
\trivia{People at Watcom were knew how to produce good compiler and they were also clever about marketing it. There was never a 1.0 or 2.0. They started directly at "version 6". This was at least one version ahead of their competitors (Borland and Microsoft). A higher number undoublty uncounciously carried a "more advanced than its competitor" feeling. Version one was also likely to features many bugs whereas the sixth installement was likely to have been battle hardened.}
\par

At a cost of \$639\footnote{Byte Magazine March 1992} (a standard price for a compiler at the time), it was not cheap but opened was a small price to pay for the door to a new world of possibilities it opened to id Software. The DOS Extender would be the corner stone bridging between PC and NeXT workstation.\\

\subsubsection{Popularity}
id Software were not the only team to value Watcom solution. Many other studios enthrusted it with their code and as a result many emblematic software of the 90s were built with Watcom technology.\\
\begin{enumerate}
\item id Softare 
       \begin{enumerate}
       \item Doom (1993)
       \item Doom II (1994)
       \end{enumerate} 
\item Blizzard Entertainment 
       \begin{enumerate}
       \item Warcraft (1994)
       \item Warcraft II (1995)
       \end{enumerate}
\item Ken Silverman's BUILD Engine based games
      \begin{enumerate}
       \item Duke Nukem 3D (1996)
       \item Shadow Warrior (1997)
       \item Blood (1997)
       \end{enumerate}
\item LucasArts Entertainment Company
      \begin{enumerate}
       \item Full Throttle (1995)
       \item The Dig (1995)
       \item Dark Forces  (1995)
       \item Rebel Assault II  (1995)    
      \end{enumerate}
\end{enumerate}
\par
